% WIP draft for Section 5 (GGRT sharpness).
\section{Proof of Theorem \ref{thm:GGRT-sharp}}

We prove that the quantitative stability estimate in Corollary~1.4 of
\cite{Gomez-Guerra-Ramos-Tilli}, which builds on the rigidity result of
\cite{Nicola-Tilli} characterizing the disk as the unique optimizer of the
first localization eigenvalue, is sharp. The proof proceeds by constructing a
family of domains which are perturbations of the disk, whose first localization
eigenvalue has the same first-order behaviour as the disk, but whose distance
to the class of disks is of order $\varepsilon$.

\begin{proof}[Proof of Theorem~\ref{thm:GGRT-sharp}]
Fix
\[
\lambda_0:=1-e^{-\pi},
\]
which is the eigenvalue of the localization operator of the unit disk $D_1$
associated with the ground state $h_0$. For $\varepsilon$ small, define
\[
P_\varepsilon(w):=1+\varepsilon w^2.
\]
By Theorem~\ref{thm:perturb}, applied with \(\lambda=\lambda_0\) and
\(P_\varepsilon(w)=1+\varepsilon w^2\), for all sufficiently small
\(|\varepsilon|\) there exists a domain \(U_\varepsilon\) such that
\begin{equation}\label{eq:functiona-P-eps-new}
\int_{U_\varepsilon} P_\varepsilon(z)e^{\pi \overline zw}e^{-\pi|z|^2}\,dA(z)
=
\lambda_0 P_\varepsilon(w),
\qquad w\in\C.
\end{equation}
Moreover, by construction, $U_\varepsilon$ is a $C^2$ perturbation of the unit
disk. Hence, after possibly shrinking the range of admissible $\varepsilon$,
its boundary can be written as
\begin{equation}\label{eq:def-u-eps-new}
\partial U_\varepsilon
=
\{(1+\varphi_\varepsilon(y))y:\ y\in \mathbb S^1\},
\end{equation}
where
\[
\varphi_\varepsilon\in C^2(\mathbb S^1),
\qquad
\|\varphi_\varepsilon\|_{C^2(\mathbb S^1)}\lesssim |\varepsilon|.
\]
The map
\[
\varepsilon\mapsto \varphi_\varepsilon
\]
may in fact be chosen \(C^1\) as a map from a neighborhood of \(0\) into
\(C^1(\mathbb S^1)\). To see this, we briefly retrace the proof of
Theorem~\ref{thm:perturb}. The construction proceeds in two steps: first, an
Isakov-type step which converts the weighted quadrature identity into a free
boundary problem of obstacle type, and second, the regularity/uniqueness step
which produces \(\varphi_\varepsilon\) via an implicit-function-theorem
argument. In both steps, the unknown is the solution of an equation of the form
\(F(\varepsilon,\varphi)=0\), where
\[
F:(-\varepsilon_0,\varepsilon_0)\times \mathcal U\to \mathcal V
\]
is a \(C^1\) map between open subsets of suitable Banach spaces, with
\(F(0,0)=0\) and \(\partial_\varphi F(0,0)\) invertible (the invertibility being
the linearization computed at the disk, which is a Fredholm operator of index
zero whose kernel is removed by the normalization fixing the centroid and the
total area at first order). Since the polynomial \(P_\varepsilon(z)=1+\varepsilon z^2\)
depends polynomially on \(\varepsilon\), the dependence of \(F\) on
\(\varepsilon\) is \(C^\infty\). The \(C^1\) version of the implicit function
theorem in Banach spaces then yields a \(C^1\) branch
\(\varepsilon\mapsto \varphi_\varepsilon\) into \(C^1(\mathbb S^1)\). We
therefore write
\[
\dot\varphi_0(y):=
\left.\partial_\varepsilon\right|_{\varepsilon=0}\varphi_\varepsilon(y).
\]

\medskip
\noindent
{\bf Step 1: first variation of the eigenvalue identity.} We differentiate \eqref{eq:functiona-P-eps-new} at $\varepsilon=0$.
Since \(P_\varepsilon\to 1\) in \(C^\infty\) and
\(\varepsilon\mapsto \varphi_\varepsilon\) is \(C^1\) into \(C^1(\mathbb S^1)\),
all terms in \eqref{eq:functiona-P-eps-new} are differentiable at
\(\varepsilon=0\). Since $U_\varepsilon$ is given by a normal graph over
$\mathbb S^1$, we apply Reynolds' transport formula. Concretely, in polar
coordinates \(z=re^{it}\), for any smooth test function \(H\) we have
\[
\int_{U_\varepsilon} H(z)\,dA(z)
=
\int_0^{2\pi}\int_0^{1+\varphi_\varepsilon(e^{it})}
H(re^{it})\,r\,dr\,dt.
\]
Differentiating in \(\varepsilon\) at \(\varepsilon=0\) and using the
fundamental theorem of calculus,
\[
\left.\frac{d}{d\varepsilon}\right|_{\varepsilon=0}
\int_{U_\varepsilon} H(z)\,dA(z)
=
\int_0^{2\pi}\dot\varphi_0(e^{it})H(e^{it})\,dt
=
\int_{\mathbb S^1}\dot\varphi_0(z)H(z)\,d\mathcal H^1(z),
\]
where the factor \(r=1\) on the unit circle has been absorbed into the
arc-length measure. This is precisely Reynolds' transport identity:
\[
\left.\frac{d}{d\varepsilon}\right|_{\varepsilon=0}
\int_{U_\varepsilon} H(z)\,dA(z)
=
\int_{\mathbb S^1} \dot\varphi_0(z)H(z)\,d\mathcal H^1(z).
\]
Applying this to
\[
H(z)=e^{\pi \overline zw}e^{-\pi|z|^2},
\]
and differentiating the factor $P_\varepsilon(z)=1+\varepsilon z^2$, we obtain
\begin{equation}\label{eq:first-variation-identity-new}
\int_{\mathbb S^1}
\dot\varphi_0(z)e^{\pi \overline zw}e^{-\pi|z|^2}\,d\mathcal H^1(z)
+
\int_{D_1} z^2 e^{\pi \overline zw}e^{-\pi|z|^2}\,dA(z)
=
\lambda_0 w^2.
\end{equation}
Since $D_1$ is radial, the monomial $z^2$ is an eigenfunction of the
localization operator of $D_1$; therefore there exists $\lambda_2>0$ such that
\[
\int_{D_1} z^2 e^{\pi \overline zw}e^{-\pi|z|^2}\,dA(z)=\lambda_2 w^2.
\]
Hence \eqref{eq:first-variation-identity-new} becomes
\begin{equation}\label{eq:boundary-identity-new}
\int_{\mathbb S^1}
\dot\varphi_0(z)e^{\pi \overline zw}e^{-\pi|z|^2}\,d\mathcal H^1(z)
=
(\lambda_0-\lambda_2)w^2.
\end{equation}
In fact, \(\lambda_2\) can be computed explicitly. Expanding
\(e^{\pi \overline zw}=\sum_{n\geq 0}\frac{\pi^n}{n!}\overline z^{\,n}w^n\)
and using orthogonality of the monomials \(\overline z^{\,n}\) over the disk,
the only term contributing to \(w^2\) is the one with \(n=2\):
\[
\int_{D_1} z^2 e^{\pi \overline zw}e^{-\pi|z|^2}\,dA(z)
=
\frac{\pi^2}{2}w^2\int_{D_1}z^2\overline z^{\,2}e^{-\pi|z|^2}\,dA(z)
=
\frac{\pi^2}{2}w^2\int_{D_1}|z|^4 e^{-\pi|z|^2}\,dA(z).
\]
Passing to polar coordinates and integrating by parts twice,
\[
\int_{D_1}|z|^4 e^{-\pi|z|^2}\,dA(z)
=
2\pi\int_0^1 r^4 e^{-\pi r^2}\,r\,dr
=
\pi\int_0^1 s^2 e^{-\pi s}\,ds
\]
where we substituted \(s=r^2\). A direct computation, integrating by parts
twice, yields
\[
\pi\int_0^1 s^2 e^{-\pi s}\,ds
=
\frac{2}{\pi^2}\Bigl(1-\Bigl(1+\pi+\frac{\pi^2}{2}\Bigr)e^{-\pi}\Bigr).
\]
Multiplying by \(\frac{\pi^2}{2}\), we obtain
\[
\lambda_2
=
\frac{\pi^2}{2}\int_{D_1}|z|^4e^{-\pi|z|^2}\,dA(z)
=
1-\Bigl(1+\pi+\frac{\pi^2}{2}\Bigr)e^{-\pi}.
\]
Therefore
\[
\lambda_0-\lambda_2
=
(1-e^{-\pi})-\Bigl(1-\Bigl(1+\pi+\frac{\pi^2}{2}\Bigr)e^{-\pi}\Bigr)
=
\pi\Bigl(1+\frac{\pi}{2}\Bigr)e^{-\pi}>0,
\]
so the coefficient of \(w^2\) on the right-hand side is indeed nonzero.
In particular, since $|z|=1$ on $\mathbb S^1$, we may rewrite this as
\begin{equation}\label{eq:boundary-identity-circle-new}
e^{-\pi}
\int_{\mathbb S^1}
\dot\varphi_0(z)e^{\pi \overline zw}\,d\mathcal H^1(z)
=
(\lambda_0-\lambda_2)w^2.
\end{equation}

Setting $w=0$ gives
\begin{equation}\label{eq:cancellation-f-0-new}
\int_{\mathbb S^1}\dot\varphi_0\,d\mathcal H^1=0.
\end{equation}

\medskip
\noindent
{\bf Step 2: identification of the leading mode.} Expanding
\[
e^{\pi \overline zw}
=
\sum_{n\geq 0}\frac{\pi^n}{n!}\overline z^{\,n}w^n,
\]
and comparing coefficients in \eqref{eq:boundary-identity-circle-new}, we find
\[
\int_{\mathbb S^1}\dot\varphi_0(z)\overline z^{\,n}\,d\mathcal H^1(z)=0
\qquad \text{for all } n\neq 2,
\]
whereas the coefficient corresponding to $n=2$ is nonzero. Since
$\dot\varphi_0$ is real-valued, it follows that
\begin{equation}\label{eq:mode-decomposition-new}
\dot\varphi_0(e^{it})
=
a\cos(2t)+b\sin(2t)
\end{equation}
for some pair $(a,b)\neq (0,0)$. In particular,
\begin{equation}\label{eq:l1-nontrivial-new}
\int_{\mathbb S^1} |\dot\varphi_0|\,d\mathcal H^1>0.
\end{equation}

\medskip
\noindent
{\bf Step 3: normalization of the area.} The area of $U_\varepsilon$ is given by
\begin{equation}\label{eq:area-u-eps-new}
|U_\varepsilon|
=
\frac12\int_{\mathbb S^1} (1+\varphi_\varepsilon(y))^2\,d\mathcal H^1(y).
\end{equation}
Since
\[
\varphi_\varepsilon
=
\varepsilon \dot\varphi_0 + O(\varepsilon^2)
\quad \text{in } C^1(\mathbb S^1),
\]
and \eqref{eq:cancellation-f-0-new} gives the vanishing of the first-order
term, we obtain from \eqref{eq:area-u-eps-new}
\begin{equation}\label{eq:measure-close-u-eps-new}
|U_\varepsilon|=\pi+O(\varepsilon^2).
\end{equation}

Let
\[
r_\varepsilon:=\sqrt{\frac{\pi}{|U_\varepsilon|}},
\qquad
\widetilde U_\varepsilon:=r_\varepsilon U_\varepsilon.
\]
Then $|\widetilde U_\varepsilon|=\pi$, and by
\eqref{eq:measure-close-u-eps-new},
\[
r_\varepsilon=1+O(\varepsilon^2).
\]
Hence $\widetilde U_\varepsilon$ still admits a normal graph representation
\[
\partial \widetilde U_\varepsilon
=
\{(1+\widetilde\varphi_\varepsilon(y))y:\ y\in \mathbb S^1\},
\]
with
\[
\widetilde\varphi_\varepsilon
=
\varepsilon \dot\varphi_0 + O(\varepsilon^2)
\quad \text{in } C^1(\mathbb S^1).
\]
Thus the area normalization does not alter the first-order deformation.

\medskip
\noindent
{\bf Step 4: distance from the centered disk.} Since both $\widetilde U_\varepsilon$ and $D_1$ are star-shaped with respect to
the origin, the symmetric difference may be computed in polar coordinates:
\[
|\widetilde U_\varepsilon \triangle D_1|
=
\frac12\int_{\mathbb S^1}
\left|(1+\widetilde\varphi_\varepsilon(y))^2-1\right|
\,d\mathcal H^1(y).
\]
Since \(\widetilde\varphi_\varepsilon=\varepsilon\dot\varphi_0+O(\varepsilon^2)\)
in \(C^1(\mathbb S^1)\), we have in particular
\(\|\widetilde\varphi_\varepsilon\|_{L^\infty(\mathbb S^1)}\lesssim |\varepsilon|\),
whence \(\widetilde\varphi_\varepsilon^2=O(\varepsilon^2)\) in
\(L^\infty(\mathbb S^1)\). Hence
\[
(1+\widetilde\varphi_\varepsilon)^2-1
=
2\widetilde\varphi_\varepsilon+\widetilde\varphi_\varepsilon^2
=
2\varepsilon \dot\varphi_0 + O(\varepsilon^2)
\quad \text{in } L^\infty(\mathbb S^1),
\]
where the implied constant depends only on the \(C^1\)-norm bound on
\(\widetilde\varphi_\varepsilon\). Therefore
\[
|\widetilde U_\varepsilon \triangle D_1|
=
\frac12\int_{\mathbb S^1}
\left|2\varepsilon \dot\varphi_0 + O(\varepsilon^2)\right|
\,d\mathcal H^1.
\]
Set \(\alpha:=\int_{\mathbb S^1}|\dot\varphi_0|\,d\mathcal H^1\). By
\eqref{eq:l1-nontrivial-new}, \(\alpha>0\). The reverse triangle inequality in
\(L^1\) and the \(O(\varepsilon^2)\) remainder give
\[
\int_{\mathbb S^1}\bigl|2\varepsilon\dot\varphi_0+O(\varepsilon^2)\bigr|\,d\mathcal H^1
\geq
2|\varepsilon|\alpha-O(\varepsilon^2)
\]
and similarly the upper bound
\[
\int_{\mathbb S^1}\bigl|2\varepsilon\dot\varphi_0+O(\varepsilon^2)\bigr|\,d\mathcal H^1
\leq
2|\varepsilon|\alpha+O(\varepsilon^2).
\]
Consequently there exist constants $c,C>0$ such that
\begin{equation}\label{eq:distance-centered-disk-new}
c|\varepsilon|
\leq
|\widetilde U_\varepsilon \triangle D_1|
\leq
C|\varepsilon|
\end{equation}
for all sufficiently small $\varepsilon$.

\medskip
\noindent
{\bf Step 5: comparison with arbitrary disks.} Let $D$ be any disk of area $\pi$. Then $D=D_1+a$ for some $a\in \C$. We first consider disks whose centers are not too close to the origin.
Since
\[
|D_1\triangle (D_1+a)|
\gtrsim \min\{|a|,1\},
\]
the triangle inequality and \eqref{eq:distance-centered-disk-new} imply that
there exists $K\gg 1$ such that, if $|a|\geq K|\varepsilon|$, then
\begin{equation}\label{eq:far-disks-new}
|\widetilde U_\varepsilon \triangle D|
\geq
|D_1\triangle D|-|\widetilde U_\varepsilon\triangle D_1|
\gtrsim
|a|-C|\varepsilon|
\gtrsim
|\varepsilon|.
\end{equation}

It remains to consider disks whose centers satisfy $|a|\leq K|\varepsilon|$.
Write
\[
a=\eta \varepsilon e^{i\theta},
\qquad
|\eta|\leq K,
\quad
\theta\in [0,2\pi).
\]
We claim that, for $\varepsilon$ sufficiently small, the translated domain
\[
\widetilde U_\varepsilon-a
\]
is still a normal graph over $\mathbb S^1$:
\[
\partial(\widetilde U_\varepsilon-a)
=
\{(1+\psi_\varepsilon^{\theta,\eta}(y))y:\ y\in \mathbb S^1\},
\]
with
\[
\psi_\varepsilon^{\theta,\eta}\in C^1(\mathbb S^1),
\qquad
\|\psi_\varepsilon^{\theta,\eta}\|_{C^1}\lesssim |\varepsilon|,
\]
uniformly in $(\theta,\eta)$ with $|\eta|\leq K$. This follows because
$\widetilde U_\varepsilon$ is already an $O(\varepsilon)$-$C^1$ perturbation of
$D_1$, and the translation vector $a$ is also of order $\varepsilon$. We now compute the first-order term of $\psi_\varepsilon^{\theta,\eta}$.
A translation by $a$ contributes, to first order in $\varepsilon$, the normal
displacement
\[
-\eta \Re(e^{-i\theta}y).
\]
Indeed, for \(y\in \mathbb S^1\),
\[
|y-a|
=
\bigl(1-2\Re(\overline y a)+|a|^2\bigr)^{1/2}
=
1-\Re(\overline y a)+O(|a|^2),
\]
uniformly in \(y\), and here \(a=\eta\varepsilon e^{i\theta}\).
Therefore
\begin{equation}\label{eq:psi-expansion-new}
\psi_\varepsilon^{\theta,\eta}(y)
=
\varepsilon\Big(\dot\varphi_0(y)-\eta \Re(e^{-i\theta}y)\Big)
+
O(\varepsilon^2)
\quad \text{in } C^1(\mathbb S^1),
\end{equation}
uniformly for $|\eta|\leq K$.

Define
\[
g^{\theta,\eta}(y)
:=
\dot\varphi_0(y)-\eta \Re(e^{-i\theta}y).
\]
Using \eqref{eq:mode-decomposition-new}, we see that $g^{\theta,\eta}$ belongs
to the finite-dimensional space spanned by the first and second harmonics.
More precisely,
\[
g^{\theta,\eta}(e^{it})
=
a\cos(2t)+b\sin(2t)-\eta \cos(t-\theta).
\]
The second-harmonic part is independent of $(\theta,\eta)$ and nonzero, because
$(a,b)\neq (0,0)$. To extract a uniform \(L^1\) lower bound, note that
\(g^{\theta,\eta}\) lies in the \(4\)-dimensional space
\[
V:=\mathrm{span}\bigl\{\cos t,\sin t,\cos 2t,\sin 2t\bigr\}\subset L^1(\mathbb S^1).
\]
On this finite-dimensional space, all norms are equivalent; in particular, the
\(L^1(\mathbb S^1)\) and \(L^\infty(\mathbb S^1)\) norms are equivalent, and
convergence in coefficients is equivalent to convergence in any of these norms. The map
\[
(\theta,\eta)\in [0,2\pi)\times [-K,K]\longmapsto g^{\theta,\eta}\in V
\]
is continuous, and its image
\[
\mathcal G_K:=\{g^{\theta,\eta}:\ \theta\in[0,2\pi),\ |\eta|\leq K\}
\]
is the continuous image of a compact set, hence compact in \(V\). Moreover, the
\((\cos 2t,\sin 2t)\)-coefficients of every \(g^{\theta,\eta}\) are exactly
\((a,b)\neq (0,0)\), independent of \((\theta,\eta)\); thus \(0\notin
\mathcal G_K\). The function \(g\mapsto\|g\|_{L^1(\mathbb S^1)}\) is continuous
and strictly positive on \(V\setminus\{0\}\), so it attains a positive minimum
on the compact set \(\mathcal G_K\). It follows that there exists
$\delta_K>0$ such that
\begin{equation}\label{eq:uniform-l1-lower-new}
\int_{\mathbb S^1}|g^{\theta,\eta}(y)|\,d\mathcal H^1(y)\geq \delta_K
\qquad
\text{for all }\theta\in[0,2\pi),\ |\eta|\leq K.
\end{equation}

Now,
\[
|\widetilde U_\varepsilon \triangle D|
=
|(\widetilde U_\varepsilon-a)\triangle D_1|.
\]
Since $\widetilde U_\varepsilon-a$ is star-shaped with radial graph
$1+\psi_\varepsilon^{\theta,\eta}$, the same polar-coordinate computation as in
Step~4 gives
\[
|(\widetilde U_\varepsilon-a)\triangle D_1|
=
\frac12\int_{\mathbb S^1}
\left|(1+\psi_\varepsilon^{\theta,\eta}(y))^2-1\right|
\,d\mathcal H^1(y).
\]
Using \eqref{eq:psi-expansion-new},
\[
(1+\psi_\varepsilon^{\theta,\eta})^2-1
=
2\varepsilon g^{\theta,\eta}+O(\varepsilon^2)
\quad \text{in } L^\infty(\mathbb S^1),
\]
uniformly for $|\eta|\leq K$. Thus, by \eqref{eq:uniform-l1-lower-new},
there exists $c_K>0$ such that
\begin{equation}\label{eq:near-disks-new}
|\widetilde U_\varepsilon \triangle D|
=
|(\widetilde U_\varepsilon-a)\triangle D_1|
\geq
c_K |\varepsilon|
\end{equation}
whenever $|a|\leq K|\varepsilon|$. Combining \eqref{eq:far-disks-new} and \eqref{eq:near-disks-new}, we conclude
that
\[
\inf_{\substack{D \text{ disk}\\ |D|=\pi}}
|\widetilde U_\varepsilon \triangle D|
\gtrsim |\varepsilon|.
\]
Since \(r_\varepsilon=1+O(\varepsilon^2)\), the dilation from \(U_\varepsilon\)
to \(\widetilde U_\varepsilon\) changes the set only by
\[
|\widetilde U_\varepsilon\triangle U_\varepsilon|=O(\varepsilon^2),
\]
and therefore does not affect the first-order distance estimate. Thus the
area-normalized family \(\widetilde U_\varepsilon\) remains at
symmetric-difference distance \(\gtrsim |\varepsilon|\) from the class of
disks.

\medskip
\noindent
{\bf Step 6: square-root sharpness of the deficit estimate.} It remains to translate the distance estimate into the corresponding deficit
estimate. Recall that, by construction via Theorem~\ref{thm:perturb}, the
first localization eigenvalue of \(U_\varepsilon\) equals \(\lambda_0\)
\emph{exactly}, since \eqref{eq:functiona-P-eps-new} states that
\(P_\varepsilon\) is an eigenfunction with eigenvalue \(\lambda_0\) for the
localization operator on \(U_\varepsilon\), and a standard variational
argument shows that this is the principal eigenvalue. By
\eqref{eq:measure-close-u-eps-new},
\(|U_\varepsilon|=\pi+O(\varepsilon^2)\), and rescaling by
\(r_\varepsilon=1+O(\varepsilon^2)\) to area \(\pi\) changes the eigenvalue
only at order \(O(\varepsilon^2)\) (since the localization operator scales
analytically in the dilation parameter, and the disk is a critical point for
this scaling). It follows that the eigenvalue deficit of
\(\widetilde U_\varepsilon\) with respect to the disk satisfies
\[
\lambda_0-\lambda(\widetilde U_\varepsilon)=O(\varepsilon^2),
\]
where \(\lambda(\widetilde U_\varepsilon)\) denotes the first localization
eigenvalue of \(\widetilde U_\varepsilon\). On the other hand, by Step~5,
the symmetric-difference asymmetry satisfies
\[
\inf_{\substack{D\text{ disk}\\ |D|=\pi}}|\widetilde U_\varepsilon\triangle D|
\gtrsim |\varepsilon|.
\]
Comparing the two,
\[
\inf_{\substack{D\text{ disk}\\ |D|=\pi}}|\widetilde U_\varepsilon\triangle D|
\gtrsim |\varepsilon|
\simeq
\sqrt{\lambda_0-\lambda(\widetilde U_\varepsilon)}.
\]
This shows that the inequality
\[
\inf_{\substack{D\text{ disk}\\ |D|=\pi}}
|\Omega\triangle D|
\lesssim
\sqrt{\lambda_0-\lambda(\Omega)}
\]
of Corollary~1.4 in \cite{Gomez-Guerra-Ramos-Tilli} is square-root-tight along
the family \(\widetilde U_\varepsilon\): the deficit is of order
\(\varepsilon^2\), the asymmetry is of order \(\varepsilon\), and so the
square root of the deficit and the asymmetry are of comparable size. This
exhibits the optimality of the exponent \(1/2\) and concludes the proof of the
sharpness of the estimate in Corollary~1.4 of
\cite{Gomez-Guerra-Ramos-Tilli}, complementing the rigidity result of
\cite{Nicola-Tilli}.

\end{proof}
